\documentclass[11pt,a4paper]{article}
\usepackage{notebooks}

\notebooktitle{DevOps Cheat Sheet}
\notebooksubtitle{Linux, Git, Docker, Kubernetes, Helm, Terraform}
\notebooktagline{Every command, grouped by the tool and by what you are trying to do.}
\notebookauthor{Bhuvnesh Verma}
\notebooksource{github.com/MasterBhuvnesh/notebooks}

\begin{document}
\notebookcover

\section{Linux}

Linux is the ground every other tool here stands on. These commands navigate the
filesystem, manage files, configure permissions and automate work from a terminal.

\subsection{Basic}

\begin{cmdtable}
pwd & Print the current working directory \\
ls & List files and directories \\
cd & Change directory \\
touch & Create an empty file \\
mkdir & Create a new directory \\
rm & Remove files or directories \\
rmdir & Remove empty directories \\
cp & Copy files or directories \\
mv & Move or rename files and directories \\
cat & Display the content of a file \\
echo & Display a line of text \\
clear & Clear the terminal screen \\
\end{cmdtable}

\subsection{Intermediate}

\begin{cmdtable}
chmod & Change file permissions \\
chown & Change file ownership \\
find & Search for files and directories \\
grep & Search for text in a file \\
wc & Count lines, words and characters in a file \\
head & Display the first few lines of a file \\
tail & Display the last few lines of a file \\
sort & Sort the contents of a file \\
uniq & Remove duplicate lines from a file \\
diff & Compare two files line by line \\
tar & Archive files into a tarball \\
zip/unzip & Compress and extract ZIP files \\
df & Display disk space usage \\
du & Display directory size \\
top & Monitor system processes in real time \\
ps & Display active processes \\
kill & Terminate a process by its PID \\
ping & Check network connectivity \\
wget & Download files from the internet \\
curl & Transfer data from or to a server \\
scp & Securely copy files between systems \\
rsync & Synchronise files and directories \\
\end{cmdtable}

\subsection{Advanced}

\begin{cmdtable}
awk & Text processing and pattern scanning \\
sed & Stream editor for filtering and transforming text \\
cut & Remove sections from each line of a file \\
tr & Translate or delete characters \\
xargs & Build and execute command lines from standard input \\
ln & Create symbolic or hard links \\
df -h & Display disk usage in human-readable format \\
free & Display memory usage \\
iostat & Display CPU and I/O statistics \\
netstat & Network statistics; \texttt{ss} is the modern alternative \\
ifconfig/ip & Configure network interfaces; \texttt{ip} is the modern alternative \\
iptables & Configure firewall rules \\
systemctl & Control the systemd system and service manager \\
journalctl & View system logs \\
crontab & Schedule recurring tasks \\
at & Schedule tasks for a specific time \\
uptime & Display system uptime \\
whoami & Display the current user \\
users & List all users currently logged in \\
hostname & Display or set the system hostname \\
env & Display environment variables \\
export & Set environment variables \\
\end{cmdtable}

\subsection{Networking}

\begin{cmdtable}
ip addr & Display or configure IP addresses \\
ip route & Show or manipulate routing tables \\
traceroute & Trace the route packets take to a host \\
nslookup & Query DNS records \\
dig & Query DNS servers \\
ssh & Connect to a remote server via SSH \\
ftp & Transfer files using the FTP protocol \\
nmap & Network scanning and discovery \\
telnet & Communicate with remote hosts \\
netcat (nc) & Read and write data over networks \\
\end{cmdtable}

\subsection{File management and search}

\begin{cmdtable}
locate & Find files quickly using a database \\
stat & Display detailed information about a file \\
tree & Display directories as a tree \\
file & Determine a file's type \\
basename & Extract the filename from a path \\
dirname & Extract the directory part of a path \\
\end{cmdtable}

\subsection{System monitoring}

\begin{cmdtable}
vmstat & Display virtual memory statistics \\
htop & Interactive process viewer, an alternative to \texttt{top} \\
lsof & List open files \\
dmesg & Print kernel ring buffer messages \\
uptime & Show how long the system has been running \\
iotop & Display real-time disk I/O by process \\
\end{cmdtable}

\subsection{Package management}

\begin{cmdtable}
apt & Package manager for Debian-based distributions \\
yum/dnf & Package manager for RHEL-based distributions \\
snap & Manage snap packages \\
rpm & Manage RPM packages \\
\end{cmdtable}

\subsection{Disk and filesystem}

\begin{cmdtable}
mount/umount & Mount or unmount filesystems \\
fsck & Check and repair filesystems \\
mkfs & Create a new filesystem \\
blkid & Display information about block devices \\
lsblk & List information about block devices \\
parted & Manage partitions interactively \\
\end{cmdtable}

\subsection{Scripting and automation}

\begin{cmdtable}
bash & Command interpreter and scripting shell \\
sh & Legacy shell interpreter \\
cron & Automate recurring tasks \\
alias & Create shortcuts for commands \\
source & Execute commands from a file in the current shell \\
\end{cmdtable}

\subsection{Development and debugging}

\begin{cmdtable}
gcc & Compile C programs \\
make & Build and manage projects \\
strace & Trace system calls and signals \\
gdb & Debug programs \\
git & Version control system \\
vim/nano & Text editors for scripting and editing \\
\end{cmdtable}

\subsection{Other useful commands}

\begin{cmdtable}
uptime & Display system uptime \\
date & Display or set the system date and time \\
cal & Display a calendar \\
man & Display the manual for a command \\
history & Show previously executed commands \\
alias & Create custom shortcuts for commands \\
\end{cmdtable}

\newpage
\section{Git}

Git tracks every change, lets a team work on the same code without collisions, and
lets you undo a mistake. These commands cover the everyday work and the recovery
tools you reach for when something has gone wrong.

\subsection{Basic}

\begin{cmdtable3}
git init & Initialises a new Git repository in the current directory & git init \\
git clone & Copies a remote repository to the local machine & git clone https://github.com/user/repo.git \\
git status & Displays the state of the working directory and staging area & git status \\
git add & Adds changes to the staging area & git add file.txt \\
git commit & Records changes to the repository & git commit -m "Initial commit" \\
git config & Configures user settings, such as name and email & git config -{}-global user.name "Your Name" \\
git log & Shows the commit history & git log \\
git show & Displays detailed information about a specific commit & git show <commit-hash> \\
git diff & Shows changes between commits, the working directory and the staging area & git diff \\
git reset & Unstages changes or resets commits & git reset HEAD file.txt \\
\end{cmdtable3}

\subsection{Branching and merging}

\begin{cmdtable3}
git branch & Lists branches or creates a new branch & git branch feature-branch \\
git checkout & Switches between branches or restores files & git checkout feature-branch \\
git switch & Switches branches, the modern alternative to \texttt{git checkout} & git switch feature-branch \\
git merge & Combines changes from one branch into another & git merge feature-branch \\
git rebase & Moves or combines commits from one branch onto another & git rebase main \\
git cherry-pick & Applies specific commits from one branch to another & git cherry-pick <commit-hash> \\
\end{cmdtable3}

\subsection{Remote repositories}

\begin{cmdtable3}
git remote & Manages remote repository connections & git remote add origin https://github.com/user/repo.git \\
git push & Sends changes to a remote repository & git push origin main \\
git pull & Fetches and merges changes from a remote repository & git pull origin main \\
git fetch & Downloads changes from a remote repository without merging & git fetch origin \\
git remote -v & Lists the URLs of remote repositories & git remote -v \\
\end{cmdtable3}

\subsection{Stashing and cleaning}

\begin{cmdtable3}
git stash & Temporarily saves changes not yet committed & git stash \\
git stash pop & Applies stashed changes and removes them from the stash list & git stash pop \\
git stash list & Lists all stashes & git stash list \\
git clean & Removes untracked files from the working directory & git clean -f \\
\end{cmdtable3}

\subsection{Tagging}

\begin{cmdtable3}
git tag & Creates a tag for a specific commit & git tag -a v1.0 -m "Version 1.0" \\
git tag -d & Deletes a tag & git tag -d v1.0 \\
git push -{}-tags & Pushes tags to a remote repository & git push origin -{}-tags \\
\end{cmdtable3}

\subsection{Advanced}

\begin{cmdtable3}
git bisect & Finds the commit that introduced a bug & git bisect start \\
git blame & Shows which commit and author modified each line of a file & git blame file.txt \\
git reflog & Shows a log of changes to the tip of branches & git reflog \\
git submodule & Manages external repositories as submodules & git submodule add https://github.com/user/repo.git \\
git archive & Creates an archive of the repository files & git archive -{}-format=zip HEAD > archive.zip \\
git gc & Cleans up unnecessary files and optimises the repository & git gc \\
\end{cmdtable3}

\subsection{GitHub CLI}

\begin{cmdtable3}
gh auth login & Logs into GitHub via the command line & gh auth login \\
gh repo clone & Clones a GitHub repository & gh repo clone user/repo \\
gh issue list & Lists issues in a GitHub repository & gh issue list \\
gh pr create & Creates a pull request on GitHub & gh pr create -{}-title "New Feature" -{}-body "Description" \\
gh repo create & Creates a new GitHub repository & gh repo create my-repo \\
\end{cmdtable3}

\newpage
\section{Docker}

Docker packages an application and everything it needs into a portable container, so
it runs the same way on a laptop, in CI and in production.

\subsection{Basic}

\begin{cmdtable3}[4.2cm]
docker -{}-version & Displays the installed Docker version & docker -{}-version \\
docker info & Shows system-wide information, such as the number of containers and images & docker info \\
docker pull & Downloads an image from a registry, Docker Hub by default & docker pull ubuntu:latest \\
docker images & Lists all downloaded images & docker images \\
docker run & Creates and starts a new container from an image & docker run -it ubuntu bash \\
docker ps & Lists running containers & docker ps \\
docker ps -a & Lists all containers, including stopped ones & docker ps -a \\
docker stop & Stops a running container & docker stop container\_name \\
docker start & Starts a stopped container & docker start container\_name \\
docker rm & Removes a container & docker rm container\_name \\
docker rmi & Removes an image & docker rmi image\_name \\
docker exec & Runs a command inside a running container & docker exec -it container\_name bash \\
\end{cmdtable3}

\subsection{Intermediate}

\begin{cmdtable3}[4.6cm]
docker build & Builds an image from a Dockerfile & docker build -t my\_image . \\
docker commit & Creates a new image from a container's changes & docker commit container\_name my\_image:tag \\
docker logs & Fetches logs from a container & docker logs container\_name \\
docker inspect & Returns detailed information about a container or image & docker inspect container\_name \\
docker stats & Displays live resource usage of running containers & docker stats \\
docker cp & Copies files between a container and the host & docker cp container\_name:/path/in/container /path/on/host \\
docker rename & Renames a container & docker rename old\_name new\_name \\
docker network ls & Lists all Docker networks & docker network ls \\
docker network create & Creates a new Docker network & docker network create my\_network \\
docker network inspect & Shows details about a Docker network & docker network inspect my\_network \\
docker network connect & Connects a container to a network & docker network connect my\_network container\_name \\
docker volume ls & Lists all Docker volumes & docker volume ls \\
docker volume create & Creates a new Docker volume & docker volume create my\_volume \\
docker volume inspect & Provides details about a volume & docker volume inspect my\_volume \\
docker volume rm & Removes a Docker volume & docker volume rm my\_volume \\
\end{cmdtable3}

\subsection{Advanced}

\begin{cmdtable3}[4.6cm]
docker-compose up & Starts services defined in a \texttt{docker-compose.yml} file & docker-compose up \\
docker-compose down & Stops and removes services defined in a \texttt{docker-compose.yml} file & docker-compose down \\
docker-compose logs & Displays logs for services managed by Docker Compose & docker-compose logs \\
docker-compose exec & Runs a command in a service's container & docker-compose exec service\_name bash \\
docker save & Exports an image to a tar file & docker save -o my\_image.tar my\_image:tag \\
docker load & Imports an image from a tar file & docker load < my\_image.tar \\
docker export & Exports a container's filesystem as a tar file & docker export container\_name > container.tar \\
docker import & Creates an image from an exported container & docker import container.tar my\_new\_image \\
docker system df & Displays disk usage by Docker objects & docker system df \\
docker system prune & Cleans up unused images, containers, volumes and networks & docker system prune \\
docker tag & Assigns a new tag to an image & docker tag old\_image\_name new\_image\_name \\
docker push & Uploads an image to a registry & docker push my\_image:tag \\
docker login & Logs into a Docker registry & docker login \\
docker logout & Logs out of a Docker registry & docker logout \\
docker swarm init & Initialises a Docker Swarm mode cluster & docker swarm init \\
docker service create & Creates a new service in Swarm mode & docker service create -{}-name my\_service nginx \\
docker stack deploy & Deploys a stack using a Compose file in Swarm mode & docker stack deploy -c docker-compose.yml my\_stack \\
docker stack rm & Removes a stack in Swarm mode & docker stack rm my\_stack \\
docker checkpoint create & Creates a checkpoint for a container & docker checkpoint create container\_name checkpoint\_name \\
docker checkpoint ls & Lists checkpoints for a container & docker checkpoint ls container\_name \\
docker checkpoint rm & Removes a checkpoint & docker checkpoint rm container\_name checkpoint\_name \\
\end{cmdtable3}

\newpage
\section{Kubernetes}

Kubernetes is the conductor of a container orchestra. It automates the deployment,
scaling and management of containerised applications across a cluster of machines.

\subsection{Basic}

\begin{cmdtable3}[4.4cm]
kubectl version & Displays the Kubernetes client and server version & kubectl version -{}-short \\
kubectl cluster-info & Shows information about the cluster & kubectl cluster-info \\
kubectl get nodes & Lists all nodes in the cluster & kubectl get nodes \\
kubectl get pods & Lists all pods in the default namespace & kubectl get pods \\
kubectl get services & Lists all services in the default namespace & kubectl get services \\
kubectl get namespaces & Lists all namespaces in the cluster & kubectl get namespaces \\
kubectl describe pod & Shows detailed information about a specific pod & kubectl describe pod pod-name \\
kubectl logs & Displays logs for a specific pod & kubectl logs pod-name \\
kubectl create namespace & Creates a new namespace & kubectl create namespace my-namespace \\
kubectl delete pod & Deletes a specific pod & kubectl delete pod pod-name \\
\end{cmdtable3}

\subsection{Intermediate}

\begin{cmdtable3}[4.4cm]
kubectl apply & Applies changes defined in a YAML file & kubectl apply -f deployment.yaml \\
kubectl delete & Deletes resources defined in a YAML file & kubectl delete -f deployment.yaml \\
kubectl scale & Scales a deployment to the desired number of replicas & kubectl scale deployment my-deployment -{}-replicas=3 \\
kubectl expose & Exposes a pod or deployment as a service & kubectl expose deployment my-deployment -{}-type=LoadBalancer -{}-port=80 \\
kubectl exec & Executes a command in a running pod & kubectl exec -it pod-name -{}- /bin/bash \\
kubectl port-forward & Forwards a local port to a port in a pod & kubectl port-forward pod-name 8080:80 \\
kubectl get configmaps & Lists all ConfigMaps in the namespace & kubectl get configmaps \\
kubectl get secrets & Lists all Secrets in the namespace & kubectl get secrets \\
kubectl edit & Edits a resource definition directly in the editor & kubectl edit deployment my-deployment \\
kubectl rollout status & Displays the status of a deployment rollout & kubectl rollout status deployment/my-deployment \\
\end{cmdtable3}

\subsection{Advanced}

\begin{cmdtable3}[4.6cm]
kubectl rollout undo & Rolls back a deployment to a previous revision & kubectl rollout undo deployment/my-deployment \\
kubectl top nodes & Shows resource usage for nodes & kubectl top nodes \\
kubectl top pods & Displays resource usage for pods & kubectl top pods \\
kubectl cordon & Marks a node as unschedulable & kubectl cordon node-name \\
kubectl uncordon & Marks a node as schedulable & kubectl uncordon node-name \\
kubectl drain & Safely evicts all pods from a node & kubectl drain node-name -{}-ignore-daemonsets \\
kubectl taint & Adds a taint to a node to control pod placement & kubectl taint nodes node-name key=value:NoSchedule \\
kubectl get events & Lists all events in the cluster & kubectl get events \\
kubectl apply -k & Applies resources from a kustomization directory & kubectl apply -k ./kustomization-dir/ \\
kubectl config view & Displays the kubeconfig file & kubectl config view \\
kubectl config use-context & Switches the active context in kubeconfig & kubectl config use-context my-cluster \\
kubectl debug & Creates a debugging session for a pod & kubectl debug pod-name \\
kubectl delete namespace & Deletes a namespace and its resources & kubectl delete namespace my-namespace \\
kubectl patch & Updates a resource using a patch & kubectl patch deployment my-deployment -p '\{"spec": \{"replicas": 2\}\}' \\
kubectl rollout history & Shows the rollout history of a deployment & kubectl rollout history deployment my-deployment \\
kubectl autoscale & Automatically scales a deployment based on resource usage & kubectl autoscale deployment my-deployment -{}-cpu-percent=50 -{}-min=1 -{}-max=10 \\
kubectl label & Adds or modifies a label on a resource & kubectl label pod pod-name environment=production \\
kubectl annotate & Adds or modifies an annotation on a resource & kubectl annotate pod pod-name description="My app pod" \\
kubectl delete pv & Deletes a PersistentVolume & kubectl delete pv my-pv \\
kubectl get ingress & Lists all Ingress resources in the namespace & kubectl get ingress \\
kubectl create configmap & Creates a ConfigMap from a file or literal values & kubectl create configmap my-config -{}-from-literal=key1=value1 \\
kubectl create secret & Creates a Secret from a file or literal values & kubectl create secret generic my-secret -{}-from-literal=password=myPassword \\
kubectl api-resources & Lists all available API resources in the cluster & kubectl api-resources \\
kubectl api-versions & Lists all API versions supported by the cluster & kubectl api-versions \\
kubectl get crds & Lists all CustomResourceDefinitions & kubectl get crds \\
\end{cmdtable3}

\newpage
\section{Helm}

Helm is the package manager for Kubernetes. It installs and manages applications
packaged as charts, the way \texttt{apt} installs packages on Debian.

\subsection{Basic}

\begin{cmdtable3}[4.4cm]
helm help & Displays help for the Helm CLI or a specific command & helm help \\
helm version & Shows the Helm client and server version & helm version \\
helm repo add & Adds a new chart repository & helm repo add stable https://charts.helm.sh/stable \\
helm repo update & Updates all chart repositories to the latest version & helm repo update \\
helm repo list & Lists all the repositories added to Helm & helm repo list \\
helm search hub & Searches for charts on Helm Hub & helm search hub nginx \\
helm search repo & Searches for charts in the repositories & helm search repo stable/nginx \\
helm show chart & Displays a chart's metadata and dependencies & helm show chart stable/nginx \\
\end{cmdtable3}

\subsection{Installing and upgrading charts}

\begin{cmdtable3}[4.6cm]
helm install & Installs a chart into a Kubernetes cluster & helm install my-release stable/nginx \\
helm upgrade & Upgrades an existing release with a new version of the chart & helm upgrade my-release stable/nginx \\
helm upgrade -{}-install & Installs a chart if it is not installed, upgrades it if it is & helm upgrade -{}-install my-release stable/nginx \\
helm uninstall & Uninstalls a release & helm uninstall my-release \\
helm list & Lists all the releases installed on the cluster & helm list \\
helm status & Displays the status of a release & helm status my-release \\
\end{cmdtable3}

\subsection{Working with charts}

\begin{cmdtable3}[5.0cm]
helm create & Creates a new Helm chart in a specified directory & helm create my-chart \\
helm lint & Lints a chart to check for common errors & helm lint ./my-chart \\
helm package & Packages a chart into a \texttt{.tgz} file & helm package ./my-chart \\
helm template & Renders the Kubernetes YAML from a chart without installing it & helm template my-release ./my-chart \\
helm dependency update & Updates the dependencies in \texttt{Chart.yaml} & helm dependency update ./my-chart \\
helm dependency build & Builds dependencies for a chart & helm dependency build ./my-chart \\
helm dependency list & Lists all dependencies for a chart & helm dependency list ./my-chart \\
\end{cmdtable3}

\subsection{Inspecting and testing}

\begin{cmdtable3}[4.4cm]
helm get all & Gets all information for a release, values and templates included & helm get all my-release \\
helm get values & Displays the values used in a release & helm get values my-release \\
helm test & Runs the tests defined in a chart & helm test my-release \\
helm template -{}-debug & Renders manifests and includes debug output & helm template my-release ./my-chart -{}-debug \\
helm install -{}-dry-run & Simulates an installation without performing it & helm install my-release stable/nginx -{}-dry-run \\
helm upgrade -{}-dry-run & Simulates an upgrade without applying it & helm upgrade my-release stable/nginx -{}-dry-run \\
\end{cmdtable3}

\subsection{Values and customisation}

\begin{cmdtable3}[4.6cm]
helm install -{}-values & Installs a chart with custom values & helm install my-release stable/nginx -{}-values values.yaml \\
helm upgrade -{}-values & Upgrades a release with custom values & helm upgrade my-release stable/nginx -{}-values values.yaml \\
helm install -{}-set & Installs a chart with a value set on the command line & helm install my-release stable/nginx -{}-set replicaCount=3 \\
helm upgrade -{}-set & Upgrades a release with a value set on the command line & helm upgrade my-release stable/nginx -{}-set replicaCount=5 \\
helm uninstall -{}-purge & Removes a release and its resources, release history included & helm uninstall my-release -{}-purge \\
\end{cmdtable3}

\subsection{Namespaces and contexts}

\begin{cmdtable3}[5.0cm]
helm list -{}-namespace & Lists releases in a specific namespace & helm list -{}-namespace kube-system \\
helm list -n & Short form of the same thing & helm list -n kube-system \\
helm install -{}-namespace & Installs a chart into a specific namespace & helm install my-release stable/nginx -{}-namespace mynamespace \\
helm upgrade -{}-namespace & Upgrades a release in a specific namespace & helm upgrade my-release stable/nginx -{}-namespace mynamespace \\
helm uninstall -{}-namespace & Uninstalls a release from a specific namespace & helm uninstall my-release -{}-namespace kube-system \\
helm install -{}-kube-context & Installs a chart into the cluster named by a kubeconfig context & helm install my-release stable/nginx -{}-kube-context my-cluster \\
helm upgrade -{}-kube-context & Upgrades a release in a specific context & helm upgrade my-release stable/nginx -{}-kube-context my-cluster \\
\end{cmdtable3}

\subsection{Repositories and publishing}

\begin{cmdtable3}[4.8cm]
helm repo remove & Removes a chart repository & helm repo remove stable \\
helm repo index & Creates or updates the index file for a chart repository & helm repo index ./charts \\
helm package -{}-sign & Packages a chart and signs it with a GPG key & helm package ./my-chart -{}-sign -{}-key my-key-id \\
helm create -{}-starter & Creates a new chart from a starter template & helm create -{}-starter https://github.com/helm/charts.git \\
helm push & Pushes a chart to a chart repository & helm push ./my-chart my-repo \\
\end{cmdtable3}

\subsection{History and rollbacks}

\begin{cmdtable3}[5.2cm]
helm rollback & Rolls a release back to a previous version & helm rollback my-release 1 \\
helm history & Displays the history of a release & helm history my-release \\
helm history -{}-max & Limits the number of versions shown in the history & helm history my-release -{}-max 5 \\
helm rollback -{}-recreate-pods & Rolls back to a previous version and recreates the pods & helm rollback my-release 2 -{}-recreate-pods \\
\end{cmdtable3}

\newpage
\section{Terraform}

Terraform builds cloud infrastructure from code. Instead of clicking through an AWS,
GCP or Azure console, servers and services are declared in configuration files and
Terraform reconciles reality with them.

\begin{cmdtable}[4.8cm]
terraform -{}-help & Display general help for the Terraform CLI \\
terraform init & Initialise the working directory and download the provider plugins the configuration needs \\
terraform validate & Check the configuration files for syntax errors and internal inconsistencies \\
terraform plan & Produce an execution plan showing what Terraform will do to reach the desired state \\
terraform apply & Apply the changes required to reach the desired state, prompting for approval first \\
terraform show & Display the state, or a saved plan, in human-readable form \\
terraform output & Display the output values defined in the configuration \\
terraform destroy & Destroy the infrastructure the configuration manages, prompting for confirmation \\
terraform refresh & Update the state file from the real infrastructure without applying any change \\
terraform taint & Mark a resource for recreation on the next apply, even if nothing about it changed \\
terraform untaint & Remove the tainted status from a resource \\
terraform state & Manage the state file: move resources between modules, remove them, inspect them \\
terraform state list & List every resource tracked in the state file \\
terraform import & Bring existing infrastructure under Terraform management \\
terraform graph & Generate a graph of the resources and their dependencies \\
terraform providers & List the providers the current configuration requires \\
\end{cmdtable}

\field{On the backend:}{The backend is configured in a \texttt{backend} block inside
\texttt{terraform \{ \}}, not from the command line. It decides where the state file lives:
locally by default, or remotely in S3, Azure Blob Storage, Terraform Cloud and others.
\texttt{terraform init} is what reads it and migrates existing state.}

\end{document}
